% Appendix fragment. Include from the main paper with:
% % Appendix fragment. Include from the main paper with:
% % Appendix fragment. Include from the main paper with:
% % Appendix fragment. Include from the main paper with:
% \input{docs/w_priorcalib_method_ablation}
% Required packages: amsmath, amssymb, bm, booktabs, graphicx

\section{SynVA-A1}
\label{Appendix:SynVA-A1}

\subsection{Model-specific data pre-processing}

SynVA-A1 is trained on the processed real aneurysm cases described in the main
text, i.e. the combined AneuX, IntrA, and CMHA cases after the common mesh and
label preprocessing. We use only those cases for which the GHD fitting step can
be computed successfully. The synthetic SynVA-P1 dataset is not used for this
model. The model-specific preprocessing has two purposes: it first converts
each aneurysm into a common GHD representation, and it then extracts the
condition variables used by the conditional VAE.

For the GHD fit, the target geometry is the aneurysm sac together with the
ostium neck. The surrounding parent vessel and adjacent branches are excluded
from the deformation target, because SynVA-A1 is intended to generate an
aneurysm anchored to a given ostium rather than a complete pathological vessel
segment. The sac mesh remains open at the ostium. The opening is not closed
for canonical-mesh construction; instead, its boundary is retained and later
used to build the ostium information required for the fitting losses.

All fits are expressed relative to one canonical aneurysm mesh. Instead of
choosing one training case as this template, we build an average canonical sac
from the training data. After an initial ostium-based alignment, rays are cast
from a common origin $\bm{c}$ along shared directions $\bm{u}_j$. If $r_{ij}$
is the intersection distance for case $i$ and ray $j$, the canonical radius is
$\bar{r}_j=\operatorname*{median}_{i} r_{ij}$ and the corresponding canonical
vertex is $\bm{v}^{\mathrm{can}}_j=\bm{c}+\bar{r}_j\bm{u}_j$. This produces a
single canonical sac with an open neck boundary and a shared graph topology for
all subsequent fits. The canonical ostium ring is extracted from this boundary.

After the canonical and target meshes have been prepared, an ostium checkpoint
is created for each case. The ostium points are mapped to nearby mesh
vertices, sorted into an ordered boundary ring, optionally smoothed in the
ostium plane, and triangulated as a simple opening patch. This patch is not
used to close the canonical sac; it provides a consistent surface object on
which the neck and opening losses can be evaluated.

Each target aneurysm is then aligned to this canonical frame before
optimization. The pre-alignment is stored as a homogeneous transformation
matrix $\bm{H}_{i}^{\mathrm{pre}}\in\mathbb{R}^{4\times4}$. It translates the
target ostium center to the canonical center, rotates the target ostium normal
onto the canonical normal, and aligns the dominant in-plane axis of the target
ostium with the corresponding canonical axis. For homogeneous coordinates
$\tilde{\bm{p}}$, the aligned point is
$\tilde{\bm{p}}_{i}^{\mathrm{align}}=\bm{H}_{i}^{\mathrm{pre}}\tilde{\bm{p}}_{i}^{\mathrm{raw}}$.
The aligned sac is saved as the fitting mesh, and the stored transformation
allows the fitted result to be related back to the processed data coordinates.
The canonical mesh and the aligned targets are normalized before the GHD
optimization.

The CVAE does not use the raw target mesh as its regression target. Instead,
the target vector is read from the GHD fitting checkpoint and consists of the
flattened GHD coefficients with $K=144$ basis functions and the fitted scale.
The fitted rotation and translation are not predicted by SynVA-A1; they are
only used to decode fitted meshes and to evaluate geometry losses in the
aligned coordinate frame.

The condition input is extracted from the same case and is centered around the
ostium. SynVA-A1 therefore does not use the complete vessel tree as input, but
only the local vessel geometry that is relevant for aneurysm attachment. To
make this vessel patch consistent with the fitted GHD representation, the
candidate vessel points are first mapped into the same local coordinate frame
as the fitted aneurysm. Distances are then measured from the ostium-ring
centroid, and a neighborhood around the ostium is retained. The neighborhood
radius is chosen from the ostium size and the local distance distribution, so
that it covers the neck region and nearby parent-vessel surface without
including the entire vessel. From this local patch, 256 points are selected by
farthest-point sampling. If the local selection contains too few points, the
full candidate surface is used as a fallback.

The remaining case information is prepared in the same local frame. The actual
network condition consists of the ordered ostium ring, compact ostium
parameters containing center, normal, radius, and eccentricity, the local
vessel point set, and a small morphology descriptor vector. We use only
selected global sac descriptors as morphology condition: aneurysm surface area
($A_A$), aneurysm volume ($V_A$), convex-hull surface area ($A_{CH}$),
convex-hull volume ($V_{CH}$), maximum diameter ($D_{\max}$), maximum height
($H_{\max}$), maximum width ($W_{\max}$), orthogonal height ($H_{\perp}$), and
orthogonal width ($W_{\perp}$). In addition, 64 ostium samples and 256 dense
samples from the segmented ostium attachment region are stored for the
geometry losses at the vessel--aneurysm interface. GHD targets and continuous
condition variables are normalized with statistics computed on the training
split.

\subsection{Architecture}

SynVA-A1 consists of a GHD fitting stage and a conditional aneurysm generator.
The fitting stage defines a compact and anatomically aligned parameterization
of each training aneurysm. The generator then learns a conditional prior over
these fitted parameters, using the local ostium, nearby vessel geometry, and
morphology descriptors as conditioning information.

\subsubsection{GHD fitting}

Our fitting procedure is based on the Graph Harmonic Deformation (GHD)
representation introduced in AneuG~\cite{ding2025two}, but adapts it from a
full aneurysm--vessel segment to the local aneurysm domain considered here.
AneuG fits a vascular structure that includes the parent artery and several
openings. In contrast, our target is the aneurysm sac together with its single
ostium neck. The parent vessel is not part of the deformation target and is
used later only as conditioning information for the generator.

We retain the core GHD representation from AneuG. Each case is fitted by
deforming the canonical sac in a shared graph-harmonic basis and by optimizing
a global similarity transform. The target of the fit is therefore not a raw
mesh with arbitrary topology, but a compact set of GHD coefficients in the
common canonical eigenspace. In our implementation this basis contains
$K=144$ graph-harmonic modes.

The general shape terms are also inherited from the original AneuG fitting
objective. We use surface Chamfer distance and surface-normal matching to align
the deformed canonical sac with the target sac. We further keep the standard
mesh regularizers used in AneuG: Laplacian smoothing, edge-length
regularization, normal-consistency regularization, and local rigidity. These
terms control the overall sac shape and discourage noisy or folded
deformations.

The main methodological change is the treatment of the ostium. Original AneuG
supervises several vessel openings and, in its full-vessel setting, can use
differentiable occupancy and differentiable centreline losses to constrain the
parent-vessel segment. Since SynVA-A1 does not fit the parent vessel, these
full-vessel losses are disabled in our sac-and-neck fitting mode. Instead, we
replace the multi-opening objective by a single-ostium objective. The ostium is
represented by an ordered boundary ring and by a reconstructed opening patch
from the ostium checkpoint. This allows us to supervise the neck directly: the
opening position is matched to the target, the opening normal is aligned using
a robust ring-based normal estimate, and an area term prevents collapse of the
neck opening. Additional overlap, boundary-smoothness, and relative-volume
terms stabilize the vessel--aneurysm interface and the global sac size.

Thus, the GHD fitting contribution in SynVA-A1 is not a new deformation
parameterization, but a new anatomical use of the AneuG fitting machinery. We
use an average canonical sac instead of a single representative template, fit
only the aneurysm sac and one open ostium, replace full-vessel centreline and
occupancy constraints by ostium-specific neck constraints, and reuse one
precomputed canonical eigenspace for all target fits.

\subsubsection{Conditional Aneurysm Generator}

SynVA-A1 keeps the compact conditional GHD-VAE backbone used in AneuG. This is
the part that remains deliberately close to the original implementation: the
network operates on GHD parameters rather than mesh vertices, uses a Gaussian
latent space with the standard reparameterization trick, and injects the
condition on the decoder side. In our case the VAE input has 433 dimensions,
corresponding to $144\times3$ GHD coefficients plus the fitted scale.

The encoder receives only this normalized GHD target. It consists of a linear
input projection, one residual MLP block, and two output heads that predict the
latent mean and log-variance. The decoder mirrors this compact design. A
108-dimensional latent sample is concatenated with the encoded condition,
passed through a linear layer and one residual MLP block, and mapped back to
the normalized GHD coefficients and scale. Both encoder and decoder use a
hidden dimension of 256. Thus, the generative backbone is still the same
AneuG-style conditional GHD-VAE: the changes are not in the basic VAE
mechanism, but in what is provided as anatomical condition and how decoded
samples are supervised.

The condition network is the main architectural extension. In the original
AneuG setup, the conditional VAE is driven by a small morphology condition.
SynVA-A1 replaces this by a local vessel--ostium condition. The local vessel
patch contains 256 points around the ostium and is encoded with a
PointNet-style branch using shared point-wise MLP layers followed by max
pooling. The compact ostium parameters, namely center, normal, radius, and
eccentricity, are encoded with a separate MLP. The ordered ostium ring is
resampled to 20 points, flattened, and encoded with another MLP so that the
decoder receives an explicit representation of the neck boundary. These
branches are fused into a 96-dimensional vessel--ostium feature, concatenated
with the selected morphology descriptors, and then projected to a
64-dimensional condition embedding before being passed to the VAE decoder.

SynVA-A1 also differs from the original AneuG model in the training signal
applied to decoded samples. The dense attachment samples from the segmented
ostium region are not used as an additional encoder branch in the final model;
instead, they supervise the generated aneurysm near the vessel--aneurysm
interface. Together with opening-center and side-of-plane losses, this forces
the predicted GHD parameters to respect the prescribed ostium geometry.
Finally, SynVA-A1 adds prior-path calibration: during training, samples drawn
from the standard normal prior are decoded under the same condition and matched
to case-level and batch-level target statistics. This trains the same sampling
path that is used at inference time, where no fitted GHD target is available.

\subsection{Training Details}

\subsubsection{GHD fitting}

GHD fitting is optimized separately for each prepared case with $K=144$ basis
functions and one ostium. We use AdamW with learning rate
$10^{-3}$ for 2000 epochs. The sac-and-neck fitting mode uses one opening and
disables the full-vessel differentiable occupancy and centreline losses from
the original AneuG fitting. The active fitting weights are summarized in
Table~\ref{tab:synva-a1-ghd-loss-weights}.

\begin{table}[t]
\centering
\caption{Loss weights used for the SynVA-A1 GHD fitting.}
\label{tab:synva-a1-ghd-loss-weights}
\scriptsize
\begin{tabular}{lclc}
\toprule
Term & Weight & Term & Weight \\
\midrule
Surface Chamfer & 20.0 & Surface normals & 16.0 \\
Laplacian & 0.1 & Edge length & 0.1 \\
Normal consistency & 350.0 & Local rigidity & 100.0 \\
Opening position & 10.0 & Opening normal & 0.1 \\
Opening area & 1.0 & Opening overlap & 1.0 \\
Opening boundary smoothness & 0.5 & Sac volume & 0.05 \\
\bottomrule
\end{tabular}
\end{table}

\subsubsection{CVAE training}

We use the same stratified and balanced split as the other models in this
work: 530 training cases, 132 validation cases, and 100 held-out test cases
with seed 42. The CVAE is trained for 4000 epochs with batch size 96 using
AdamW with learning rate $10^{-4}$ and weight decay $10^{-4}$. The latent
dimension is 108, the hidden dimension is 256, the condition embedding
dimension is 64, and the vessel-condition embedding dimension is 96. We use a
KL warm-up of 1000 epochs, free bits of 0.02, validation every 50 epochs, and
select the checkpoint with the lowest validation total loss. The selected
SynVA-A1 checkpoint is from epoch 2550.

The CVAE loss combines three parts. The posterior reconstruction path matches
the fitted GHD coefficients, scale, and decoded vertices while regularizing
the posterior with a KL term. Geometry-aware losses supervise the generated
aneurysm near the ostium through attachment-sample distance, sac-center error,
wrong-side placement relative to the ostium plane, and opening-center error.
Finally, prior-path calibration decodes direct prior samples
$\bm{z}_{\mathrm{prior}}\sim\mathcal{N}(\bm{0},\bm{I})$ during training and
matches them to case-level and batch-level target statistics. We write the
overall objective as
$\mathcal{L}_{\mathrm{CVAE}}=\mathcal{L}_{\mathrm{rec}}+\mathcal{L}_{\mathrm{geo}}+\mathcal{L}_{\mathrm{prior}}$.
The corresponding weights are listed in
Table~\ref{tab:synva-a1-cvae-loss-weights}.

\begin{table}[t]
\centering
\caption{CVAE loss weights used for SynVA-A1.}
\label{tab:synva-a1-cvae-loss-weights}
\scriptsize
\begin{tabular}{lclc}
\toprule
Term & Weight & Term & Weight \\
\midrule
Coeff. MSE & 1.0 & Attachment samples & 5.0 \\
Scale Huber & 1.0 & Sac center & 10.0 \\
Vertex MSE & 250.0 & Side & 2.0 \\
KL & 0.01 & Opening center & 10.0 \\
Prior Huber & 0.25 & Prior mean/std & 0.25 / 8.0 \\
\bottomrule
\end{tabular}
\end{table}

\subsection{Experimental setting}

All ablations are trained and evaluated with the same splits and the same
100-case test set. AneuG-Base denotes the compact GHD-VAE baseline. AneuG-Cond
adds local ostium and vessel conditioning, AneuG-Cond-Morph additionally uses
the morphology vector, and AneuG-Cond-Prior adds prior-path calibration without
morphology conditioning. SynVA-A1 is the selected model with local
conditioning, morphology conditioning, geometry-aware ostium supervision, and
prior-path calibration. SynVA-A1-Ostium uses the same model but increases the
ostium-related supervision weights.

Evaluation uses paired surface Chamfer distance with 10k sampled surface
points, mean absolute errors of morphology descriptors, and distribution-level
metrics in standardized morphology space. Lower values are better for all
reported metrics except coverage.

\subsection{Additional results}

\begin{table}[t]
\centering
\caption{Summary of held-out SynVA-A1 ablation results. Paired metrics are computed against the ground-truth sac; distribution metrics are computed in standardized morphology space. Lower is better except for coverage.}
\label{tab:synva-a1-results-summary}
\scriptsize
\setlength{\tabcolsep}{3.0pt}
\begin{tabular}{lrrrrrrrr}
\toprule
Variant & Chamfer & $D_{\max}$ & $H_{\max}$ & $W_{\max}$ & paired $\bm{z}$-L2 & coverage $\uparrow$ & MMD & Frechet \\
\midrule
AneuG-Base & 0.094099 & 0.113297 & 0.119853 & 0.074605 & 1.284607 & 0.84 & 0.006438 & 1.092620 \\
AneuG-Cond & 0.089148 & 0.106784 & 0.108650 & 0.068531 & 1.217317 & 0.81 & 0.009033 & 1.111498 \\
AneuG-Cond-Morph & \textbf{0.082113} & 0.107940 & 0.103746 & 0.061546 & 1.127876 & 0.83 & 0.010141 & 1.254649 \\
AneuG-Cond-Prior & 0.087192 & \textbf{0.084195} & 0.081362 & 0.059200 & 1.008820 & \textbf{0.87} & \textbf{0.005170} & \textbf{0.984928} \\
SynVA-A1 & 0.086362 & 0.084209 & \textbf{0.079572} & 0.057382 & \textbf{0.980470} & 0.85 & 0.005519 & 1.017597 \\
SynVA-A1-Ostium & 0.086228 & 0.088093 & 0.082618 & \textbf{0.056933} & 0.996746 & 0.83 & 0.006162 & 1.061131 \\
\bottomrule
\end{tabular}
\end{table}

The ablation shows that the individual additions affect different aspects of
the generated aneurysms. Local conditioning improves paired reconstruction,
morphology conditioning gives the lowest Chamfer distance, and prior-path
calibration is most important for distribution-level behavior. SynVA-A1 gives
the best paired morphology distance and the lowest height error, while
AneuG-Cond-Prior achieves the strongest coverage, MMD, and Frechet distance.

\IfFileExists{images/fail_cases_synva_a1.png}{%
\begin{figure}[ht]
  \centering
  \includegraphics[width=0.8\textwidth]{images/fail_cases_synva_a1.png}
  \vspace{-0.2cm}
  \caption{Representative SynVA-A1 failure cases.}
  \label{fig:synva-a1-failure-cases}
\end{figure}
}{}

% Required packages: amsmath, amssymb, bm, booktabs, graphicx

\section{SynVA-A1}
\label{Appendix:SynVA-A1}

\subsection{Model-specific data pre-processing}

SynVA-A1 is trained on the processed real aneurysm cases described in the main
text, i.e. the combined AneuX, IntrA, and CMHA cases after the common mesh and
label preprocessing. We use only those cases for which the GHD fitting step can
be computed successfully. The synthetic SynVA-P1 dataset is not used for this
model. The model-specific preprocessing has two purposes: it first converts
each aneurysm into a common GHD representation, and it then extracts the
condition variables used by the conditional VAE.

For the GHD fit, the target geometry is the aneurysm sac together with the
ostium neck. The surrounding parent vessel and adjacent branches are excluded
from the deformation target, because SynVA-A1 is intended to generate an
aneurysm anchored to a given ostium rather than a complete pathological vessel
segment. The sac mesh remains open at the ostium. The opening is not closed
for canonical-mesh construction; instead, its boundary is retained and later
used to build the ostium information required for the fitting losses.

All fits are expressed relative to one canonical aneurysm mesh. Instead of
choosing one training case as this template, we build an average canonical sac
from the training data. After an initial ostium-based alignment, rays are cast
from a common origin $\bm{c}$ along shared directions $\bm{u}_j$. If $r_{ij}$
is the intersection distance for case $i$ and ray $j$, the canonical radius is
$\bar{r}_j=\operatorname*{median}_{i} r_{ij}$ and the corresponding canonical
vertex is $\bm{v}^{\mathrm{can}}_j=\bm{c}+\bar{r}_j\bm{u}_j$. This produces a
single canonical sac with an open neck boundary and a shared graph topology for
all subsequent fits. The canonical ostium ring is extracted from this boundary.

After the canonical and target meshes have been prepared, an ostium checkpoint
is created for each case. The ostium points are mapped to nearby mesh
vertices, sorted into an ordered boundary ring, optionally smoothed in the
ostium plane, and triangulated as a simple opening patch. This patch is not
used to close the canonical sac; it provides a consistent surface object on
which the neck and opening losses can be evaluated.

Each target aneurysm is then aligned to this canonical frame before
optimization. The pre-alignment is stored as a homogeneous transformation
matrix $\bm{H}_{i}^{\mathrm{pre}}\in\mathbb{R}^{4\times4}$. It translates the
target ostium center to the canonical center, rotates the target ostium normal
onto the canonical normal, and aligns the dominant in-plane axis of the target
ostium with the corresponding canonical axis. For homogeneous coordinates
$\tilde{\bm{p}}$, the aligned point is
$\tilde{\bm{p}}_{i}^{\mathrm{align}}=\bm{H}_{i}^{\mathrm{pre}}\tilde{\bm{p}}_{i}^{\mathrm{raw}}$.
The aligned sac is saved as the fitting mesh, and the stored transformation
allows the fitted result to be related back to the processed data coordinates.
The canonical mesh and the aligned targets are normalized before the GHD
optimization.

The CVAE does not use the raw target mesh as its regression target. Instead,
the target vector is read from the GHD fitting checkpoint and consists of the
flattened GHD coefficients with $K=144$ basis functions and the fitted scale.
The fitted rotation and translation are not predicted by SynVA-A1; they are
only used to decode fitted meshes and to evaluate geometry losses in the
aligned coordinate frame.

The condition input is extracted from the same case and is centered around the
ostium. SynVA-A1 therefore does not use the complete vessel tree as input, but
only the local vessel geometry that is relevant for aneurysm attachment. To
make this vessel patch consistent with the fitted GHD representation, the
candidate vessel points are first mapped into the same local coordinate frame
as the fitted aneurysm. Distances are then measured from the ostium-ring
centroid, and a neighborhood around the ostium is retained. The neighborhood
radius is chosen from the ostium size and the local distance distribution, so
that it covers the neck region and nearby parent-vessel surface without
including the entire vessel. From this local patch, 256 points are selected by
farthest-point sampling. If the local selection contains too few points, the
full candidate surface is used as a fallback.

The remaining case information is prepared in the same local frame. The actual
network condition consists of the ordered ostium ring, compact ostium
parameters containing center, normal, radius, and eccentricity, the local
vessel point set, and a small morphology descriptor vector. We use only
selected global sac descriptors as morphology condition: aneurysm surface area
($A_A$), aneurysm volume ($V_A$), convex-hull surface area ($A_{CH}$),
convex-hull volume ($V_{CH}$), maximum diameter ($D_{\max}$), maximum height
($H_{\max}$), maximum width ($W_{\max}$), orthogonal height ($H_{\perp}$), and
orthogonal width ($W_{\perp}$). In addition, 64 ostium samples and 256 dense
samples from the segmented ostium attachment region are stored for the
geometry losses at the vessel--aneurysm interface. GHD targets and continuous
condition variables are normalized with statistics computed on the training
split.

\subsection{Architecture}

SynVA-A1 consists of a GHD fitting stage and a conditional aneurysm generator.
The fitting stage defines a compact and anatomically aligned parameterization
of each training aneurysm. The generator then learns a conditional prior over
these fitted parameters, using the local ostium, nearby vessel geometry, and
morphology descriptors as conditioning information.

\subsubsection{GHD fitting}

Our fitting procedure is based on the Graph Harmonic Deformation (GHD)
representation introduced in AneuG~\cite{ding2025two}, but adapts it from a
full aneurysm--vessel segment to the local aneurysm domain considered here.
AneuG fits a vascular structure that includes the parent artery and several
openings. In contrast, our target is the aneurysm sac together with its single
ostium neck. The parent vessel is not part of the deformation target and is
used later only as conditioning information for the generator.

We retain the core GHD representation from AneuG. Each case is fitted by
deforming the canonical sac in a shared graph-harmonic basis and by optimizing
a global similarity transform. The target of the fit is therefore not a raw
mesh with arbitrary topology, but a compact set of GHD coefficients in the
common canonical eigenspace. In our implementation this basis contains
$K=144$ graph-harmonic modes.

The general shape terms are also inherited from the original AneuG fitting
objective. We use surface Chamfer distance and surface-normal matching to align
the deformed canonical sac with the target sac. We further keep the standard
mesh regularizers used in AneuG: Laplacian smoothing, edge-length
regularization, normal-consistency regularization, and local rigidity. These
terms control the overall sac shape and discourage noisy or folded
deformations.

The main methodological change is the treatment of the ostium. Original AneuG
supervises several vessel openings and, in its full-vessel setting, can use
differentiable occupancy and differentiable centreline losses to constrain the
parent-vessel segment. Since SynVA-A1 does not fit the parent vessel, these
full-vessel losses are disabled in our sac-and-neck fitting mode. Instead, we
replace the multi-opening objective by a single-ostium objective. The ostium is
represented by an ordered boundary ring and by a reconstructed opening patch
from the ostium checkpoint. This allows us to supervise the neck directly: the
opening position is matched to the target, the opening normal is aligned using
a robust ring-based normal estimate, and an area term prevents collapse of the
neck opening. Additional overlap, boundary-smoothness, and relative-volume
terms stabilize the vessel--aneurysm interface and the global sac size.

Thus, the GHD fitting contribution in SynVA-A1 is not a new deformation
parameterization, but a new anatomical use of the AneuG fitting machinery. We
use an average canonical sac instead of a single representative template, fit
only the aneurysm sac and one open ostium, replace full-vessel centreline and
occupancy constraints by ostium-specific neck constraints, and reuse one
precomputed canonical eigenspace for all target fits.

\subsubsection{Conditional Aneurysm Generator}

SynVA-A1 keeps the compact conditional GHD-VAE backbone used in AneuG. This is
the part that remains deliberately close to the original implementation: the
network operates on GHD parameters rather than mesh vertices, uses a Gaussian
latent space with the standard reparameterization trick, and injects the
condition on the decoder side. In our case the VAE input has 433 dimensions,
corresponding to $144\times3$ GHD coefficients plus the fitted scale.

The encoder receives only this normalized GHD target. It consists of a linear
input projection, one residual MLP block, and two output heads that predict the
latent mean and log-variance. The decoder mirrors this compact design. A
108-dimensional latent sample is concatenated with the encoded condition,
passed through a linear layer and one residual MLP block, and mapped back to
the normalized GHD coefficients and scale. Both encoder and decoder use a
hidden dimension of 256. Thus, the generative backbone is still the same
AneuG-style conditional GHD-VAE: the changes are not in the basic VAE
mechanism, but in what is provided as anatomical condition and how decoded
samples are supervised.

The condition network is the main architectural extension. In the original
AneuG setup, the conditional VAE is driven by a small morphology condition.
SynVA-A1 replaces this by a local vessel--ostium condition. The local vessel
patch contains 256 points around the ostium and is encoded with a
PointNet-style branch using shared point-wise MLP layers followed by max
pooling. The compact ostium parameters, namely center, normal, radius, and
eccentricity, are encoded with a separate MLP. The ordered ostium ring is
resampled to 20 points, flattened, and encoded with another MLP so that the
decoder receives an explicit representation of the neck boundary. These
branches are fused into a 96-dimensional vessel--ostium feature, concatenated
with the selected morphology descriptors, and then projected to a
64-dimensional condition embedding before being passed to the VAE decoder.

SynVA-A1 also differs from the original AneuG model in the training signal
applied to decoded samples. The dense attachment samples from the segmented
ostium region are not used as an additional encoder branch in the final model;
instead, they supervise the generated aneurysm near the vessel--aneurysm
interface. Together with opening-center and side-of-plane losses, this forces
the predicted GHD parameters to respect the prescribed ostium geometry.
Finally, SynVA-A1 adds prior-path calibration: during training, samples drawn
from the standard normal prior are decoded under the same condition and matched
to case-level and batch-level target statistics. This trains the same sampling
path that is used at inference time, where no fitted GHD target is available.

\subsection{Training Details}

\subsubsection{GHD fitting}

GHD fitting is optimized separately for each prepared case with $K=144$ basis
functions and one ostium. We use AdamW with learning rate
$10^{-3}$ for 2000 epochs. The sac-and-neck fitting mode uses one opening and
disables the full-vessel differentiable occupancy and centreline losses from
the original AneuG fitting. The active fitting weights are summarized in
Table~\ref{tab:synva-a1-ghd-loss-weights}.

\begin{table}[t]
\centering
\caption{Loss weights used for the SynVA-A1 GHD fitting.}
\label{tab:synva-a1-ghd-loss-weights}
\scriptsize
\begin{tabular}{lclc}
\toprule
Term & Weight & Term & Weight \\
\midrule
Surface Chamfer & 20.0 & Surface normals & 16.0 \\
Laplacian & 0.1 & Edge length & 0.1 \\
Normal consistency & 350.0 & Local rigidity & 100.0 \\
Opening position & 10.0 & Opening normal & 0.1 \\
Opening area & 1.0 & Opening overlap & 1.0 \\
Opening boundary smoothness & 0.5 & Sac volume & 0.05 \\
\bottomrule
\end{tabular}
\end{table}

\subsubsection{CVAE training}

We use the same stratified and balanced split as the other models in this
work: 530 training cases, 132 validation cases, and 100 held-out test cases
with seed 42. The CVAE is trained for 4000 epochs with batch size 96 using
AdamW with learning rate $10^{-4}$ and weight decay $10^{-4}$. The latent
dimension is 108, the hidden dimension is 256, the condition embedding
dimension is 64, and the vessel-condition embedding dimension is 96. We use a
KL warm-up of 1000 epochs, free bits of 0.02, validation every 50 epochs, and
select the checkpoint with the lowest validation total loss. The selected
SynVA-A1 checkpoint is from epoch 2550.

The CVAE loss combines three parts. The posterior reconstruction path matches
the fitted GHD coefficients, scale, and decoded vertices while regularizing
the posterior with a KL term. Geometry-aware losses supervise the generated
aneurysm near the ostium through attachment-sample distance, sac-center error,
wrong-side placement relative to the ostium plane, and opening-center error.
Finally, prior-path calibration decodes direct prior samples
$\bm{z}_{\mathrm{prior}}\sim\mathcal{N}(\bm{0},\bm{I})$ during training and
matches them to case-level and batch-level target statistics. We write the
overall objective as
$\mathcal{L}_{\mathrm{CVAE}}=\mathcal{L}_{\mathrm{rec}}+\mathcal{L}_{\mathrm{geo}}+\mathcal{L}_{\mathrm{prior}}$.
The corresponding weights are listed in
Table~\ref{tab:synva-a1-cvae-loss-weights}.

\begin{table}[t]
\centering
\caption{CVAE loss weights used for SynVA-A1.}
\label{tab:synva-a1-cvae-loss-weights}
\scriptsize
\begin{tabular}{lclc}
\toprule
Term & Weight & Term & Weight \\
\midrule
Coeff. MSE & 1.0 & Attachment samples & 5.0 \\
Scale Huber & 1.0 & Sac center & 10.0 \\
Vertex MSE & 250.0 & Side & 2.0 \\
KL & 0.01 & Opening center & 10.0 \\
Prior Huber & 0.25 & Prior mean/std & 0.25 / 8.0 \\
\bottomrule
\end{tabular}
\end{table}

\subsection{Experimental setting}

All ablations are trained and evaluated with the same splits and the same
100-case test set. AneuG-Base denotes the compact GHD-VAE baseline. AneuG-Cond
adds local ostium and vessel conditioning, AneuG-Cond-Morph additionally uses
the morphology vector, and AneuG-Cond-Prior adds prior-path calibration without
morphology conditioning. SynVA-A1 is the selected model with local
conditioning, morphology conditioning, geometry-aware ostium supervision, and
prior-path calibration. SynVA-A1-Ostium uses the same model but increases the
ostium-related supervision weights.

Evaluation uses paired surface Chamfer distance with 10k sampled surface
points, mean absolute errors of morphology descriptors, and distribution-level
metrics in standardized morphology space. Lower values are better for all
reported metrics except coverage.

\subsection{Additional results}

\begin{table}[t]
\centering
\caption{Summary of held-out SynVA-A1 ablation results. Paired metrics are computed against the ground-truth sac; distribution metrics are computed in standardized morphology space. Lower is better except for coverage.}
\label{tab:synva-a1-results-summary}
\scriptsize
\setlength{\tabcolsep}{3.0pt}
\begin{tabular}{lrrrrrrrr}
\toprule
Variant & Chamfer & $D_{\max}$ & $H_{\max}$ & $W_{\max}$ & paired $\bm{z}$-L2 & coverage $\uparrow$ & MMD & Frechet \\
\midrule
AneuG-Base & 0.094099 & 0.113297 & 0.119853 & 0.074605 & 1.284607 & 0.84 & 0.006438 & 1.092620 \\
AneuG-Cond & 0.089148 & 0.106784 & 0.108650 & 0.068531 & 1.217317 & 0.81 & 0.009033 & 1.111498 \\
AneuG-Cond-Morph & \textbf{0.082113} & 0.107940 & 0.103746 & 0.061546 & 1.127876 & 0.83 & 0.010141 & 1.254649 \\
AneuG-Cond-Prior & 0.087192 & \textbf{0.084195} & 0.081362 & 0.059200 & 1.008820 & \textbf{0.87} & \textbf{0.005170} & \textbf{0.984928} \\
SynVA-A1 & 0.086362 & 0.084209 & \textbf{0.079572} & 0.057382 & \textbf{0.980470} & 0.85 & 0.005519 & 1.017597 \\
SynVA-A1-Ostium & 0.086228 & 0.088093 & 0.082618 & \textbf{0.056933} & 0.996746 & 0.83 & 0.006162 & 1.061131 \\
\bottomrule
\end{tabular}
\end{table}

The ablation shows that the individual additions affect different aspects of
the generated aneurysms. Local conditioning improves paired reconstruction,
morphology conditioning gives the lowest Chamfer distance, and prior-path
calibration is most important for distribution-level behavior. SynVA-A1 gives
the best paired morphology distance and the lowest height error, while
AneuG-Cond-Prior achieves the strongest coverage, MMD, and Frechet distance.

\IfFileExists{images/fail_cases_synva_a1.png}{%
\begin{figure}[ht]
  \centering
  \includegraphics[width=0.8\textwidth]{images/fail_cases_synva_a1.png}
  \vspace{-0.2cm}
  \caption{Representative SynVA-A1 failure cases.}
  \label{fig:synva-a1-failure-cases}
\end{figure}
}{}

% Required packages: amsmath, amssymb, bm, booktabs, graphicx

\section{SynVA-A1}
\label{Appendix:SynVA-A1}

\subsection{Model-specific data pre-processing}

SynVA-A1 is trained on the processed real aneurysm cases described in the main
text, i.e. the combined AneuX, IntrA, and CMHA cases after the common mesh and
label preprocessing. We use only those cases for which the GHD fitting step can
be computed successfully. The synthetic SynVA-P1 dataset is not used for this
model. The model-specific preprocessing has two purposes: it first converts
each aneurysm into a common GHD representation, and it then extracts the
condition variables used by the conditional VAE.

For the GHD fit, the target geometry is the aneurysm sac together with the
ostium neck. The surrounding parent vessel and adjacent branches are excluded
from the deformation target, because SynVA-A1 is intended to generate an
aneurysm anchored to a given ostium rather than a complete pathological vessel
segment. The sac mesh remains open at the ostium. The opening is not closed
for canonical-mesh construction; instead, its boundary is retained and later
used to build the ostium information required for the fitting losses.

All fits are expressed relative to one canonical aneurysm mesh. Instead of
choosing one training case as this template, we build an average canonical sac
from the training data. After an initial ostium-based alignment, rays are cast
from a common origin $\bm{c}$ along shared directions $\bm{u}_j$. If $r_{ij}$
is the intersection distance for case $i$ and ray $j$, the canonical radius is
$\bar{r}_j=\operatorname*{median}_{i} r_{ij}$ and the corresponding canonical
vertex is $\bm{v}^{\mathrm{can}}_j=\bm{c}+\bar{r}_j\bm{u}_j$. This produces a
single canonical sac with an open neck boundary and a shared graph topology for
all subsequent fits. The canonical ostium ring is extracted from this boundary.

After the canonical and target meshes have been prepared, an ostium checkpoint
is created for each case. The ostium points are mapped to nearby mesh
vertices, sorted into an ordered boundary ring, optionally smoothed in the
ostium plane, and triangulated as a simple opening patch. This patch is not
used to close the canonical sac; it provides a consistent surface object on
which the neck and opening losses can be evaluated.

Each target aneurysm is then aligned to this canonical frame before
optimization. The pre-alignment is stored as a homogeneous transformation
matrix $\bm{H}_{i}^{\mathrm{pre}}\in\mathbb{R}^{4\times4}$. It translates the
target ostium center to the canonical center, rotates the target ostium normal
onto the canonical normal, and aligns the dominant in-plane axis of the target
ostium with the corresponding canonical axis. For homogeneous coordinates
$\tilde{\bm{p}}$, the aligned point is
$\tilde{\bm{p}}_{i}^{\mathrm{align}}=\bm{H}_{i}^{\mathrm{pre}}\tilde{\bm{p}}_{i}^{\mathrm{raw}}$.
The aligned sac is saved as the fitting mesh, and the stored transformation
allows the fitted result to be related back to the processed data coordinates.
The canonical mesh and the aligned targets are normalized before the GHD
optimization.

The CVAE does not use the raw target mesh as its regression target. Instead,
the target vector is read from the GHD fitting checkpoint and consists of the
flattened GHD coefficients with $K=144$ basis functions and the fitted scale.
The fitted rotation and translation are not predicted by SynVA-A1; they are
only used to decode fitted meshes and to evaluate geometry losses in the
aligned coordinate frame.

The condition input is extracted from the same case and is centered around the
ostium. SynVA-A1 therefore does not use the complete vessel tree as input, but
only the local vessel geometry that is relevant for aneurysm attachment. To
make this vessel patch consistent with the fitted GHD representation, the
candidate vessel points are first mapped into the same local coordinate frame
as the fitted aneurysm. Distances are then measured from the ostium-ring
centroid, and a neighborhood around the ostium is retained. The neighborhood
radius is chosen from the ostium size and the local distance distribution, so
that it covers the neck region and nearby parent-vessel surface without
including the entire vessel. From this local patch, 256 points are selected by
farthest-point sampling. If the local selection contains too few points, the
full candidate surface is used as a fallback.

The remaining case information is prepared in the same local frame. The actual
network condition consists of the ordered ostium ring, compact ostium
parameters containing center, normal, radius, and eccentricity, the local
vessel point set, and a small morphology descriptor vector. We use only
selected global sac descriptors as morphology condition: aneurysm surface area
($A_A$), aneurysm volume ($V_A$), convex-hull surface area ($A_{CH}$),
convex-hull volume ($V_{CH}$), maximum diameter ($D_{\max}$), maximum height
($H_{\max}$), maximum width ($W_{\max}$), orthogonal height ($H_{\perp}$), and
orthogonal width ($W_{\perp}$). In addition, 64 ostium samples and 256 dense
samples from the segmented ostium attachment region are stored for the
geometry losses at the vessel--aneurysm interface. GHD targets and continuous
condition variables are normalized with statistics computed on the training
split.

\subsection{Architecture}

SynVA-A1 consists of a GHD fitting stage and a conditional aneurysm generator.
The fitting stage defines a compact and anatomically aligned parameterization
of each training aneurysm. The generator then learns a conditional prior over
these fitted parameters, using the local ostium, nearby vessel geometry, and
morphology descriptors as conditioning information.

\subsubsection{GHD fitting}

Our fitting procedure is based on the Graph Harmonic Deformation (GHD)
representation introduced in AneuG~\cite{ding2025two}, but adapts it from a
full aneurysm--vessel segment to the local aneurysm domain considered here.
AneuG fits a vascular structure that includes the parent artery and several
openings. In contrast, our target is the aneurysm sac together with its single
ostium neck. The parent vessel is not part of the deformation target and is
used later only as conditioning information for the generator.

We retain the core GHD representation from AneuG. Each case is fitted by
deforming the canonical sac in a shared graph-harmonic basis and by optimizing
a global similarity transform. The target of the fit is therefore not a raw
mesh with arbitrary topology, but a compact set of GHD coefficients in the
common canonical eigenspace. In our implementation this basis contains
$K=144$ graph-harmonic modes.

The general shape terms are also inherited from the original AneuG fitting
objective. We use surface Chamfer distance and surface-normal matching to align
the deformed canonical sac with the target sac. We further keep the standard
mesh regularizers used in AneuG: Laplacian smoothing, edge-length
regularization, normal-consistency regularization, and local rigidity. These
terms control the overall sac shape and discourage noisy or folded
deformations.

The main methodological change is the treatment of the ostium. Original AneuG
supervises several vessel openings and, in its full-vessel setting, can use
differentiable occupancy and differentiable centreline losses to constrain the
parent-vessel segment. Since SynVA-A1 does not fit the parent vessel, these
full-vessel losses are disabled in our sac-and-neck fitting mode. Instead, we
replace the multi-opening objective by a single-ostium objective. The ostium is
represented by an ordered boundary ring and by a reconstructed opening patch
from the ostium checkpoint. This allows us to supervise the neck directly: the
opening position is matched to the target, the opening normal is aligned using
a robust ring-based normal estimate, and an area term prevents collapse of the
neck opening. Additional overlap, boundary-smoothness, and relative-volume
terms stabilize the vessel--aneurysm interface and the global sac size.

Thus, the GHD fitting contribution in SynVA-A1 is not a new deformation
parameterization, but a new anatomical use of the AneuG fitting machinery. We
use an average canonical sac instead of a single representative template, fit
only the aneurysm sac and one open ostium, replace full-vessel centreline and
occupancy constraints by ostium-specific neck constraints, and reuse one
precomputed canonical eigenspace for all target fits.

\subsubsection{Conditional Aneurysm Generator}

SynVA-A1 keeps the compact conditional GHD-VAE backbone used in AneuG. This is
the part that remains deliberately close to the original implementation: the
network operates on GHD parameters rather than mesh vertices, uses a Gaussian
latent space with the standard reparameterization trick, and injects the
condition on the decoder side. In our case the VAE input has 433 dimensions,
corresponding to $144\times3$ GHD coefficients plus the fitted scale.

The encoder receives only this normalized GHD target. It consists of a linear
input projection, one residual MLP block, and two output heads that predict the
latent mean and log-variance. The decoder mirrors this compact design. A
108-dimensional latent sample is concatenated with the encoded condition,
passed through a linear layer and one residual MLP block, and mapped back to
the normalized GHD coefficients and scale. Both encoder and decoder use a
hidden dimension of 256. Thus, the generative backbone is still the same
AneuG-style conditional GHD-VAE: the changes are not in the basic VAE
mechanism, but in what is provided as anatomical condition and how decoded
samples are supervised.

The condition network is the main architectural extension. In the original
AneuG setup, the conditional VAE is driven by a small morphology condition.
SynVA-A1 replaces this by a local vessel--ostium condition. The local vessel
patch contains 256 points around the ostium and is encoded with a
PointNet-style branch using shared point-wise MLP layers followed by max
pooling. The compact ostium parameters, namely center, normal, radius, and
eccentricity, are encoded with a separate MLP. The ordered ostium ring is
resampled to 20 points, flattened, and encoded with another MLP so that the
decoder receives an explicit representation of the neck boundary. These
branches are fused into a 96-dimensional vessel--ostium feature, concatenated
with the selected morphology descriptors, and then projected to a
64-dimensional condition embedding before being passed to the VAE decoder.

SynVA-A1 also differs from the original AneuG model in the training signal
applied to decoded samples. The dense attachment samples from the segmented
ostium region are not used as an additional encoder branch in the final model;
instead, they supervise the generated aneurysm near the vessel--aneurysm
interface. Together with opening-center and side-of-plane losses, this forces
the predicted GHD parameters to respect the prescribed ostium geometry.
Finally, SynVA-A1 adds prior-path calibration: during training, samples drawn
from the standard normal prior are decoded under the same condition and matched
to case-level and batch-level target statistics. This trains the same sampling
path that is used at inference time, where no fitted GHD target is available.

\subsection{Training Details}

\subsubsection{GHD fitting}

GHD fitting is optimized separately for each prepared case with $K=144$ basis
functions and one ostium. We use AdamW with learning rate
$10^{-3}$ for 2000 epochs. The sac-and-neck fitting mode uses one opening and
disables the full-vessel differentiable occupancy and centreline losses from
the original AneuG fitting. The active fitting weights are summarized in
Table~\ref{tab:synva-a1-ghd-loss-weights}.

\begin{table}[t]
\centering
\caption{Loss weights used for the SynVA-A1 GHD fitting.}
\label{tab:synva-a1-ghd-loss-weights}
\scriptsize
\begin{tabular}{lclc}
\toprule
Term & Weight & Term & Weight \\
\midrule
Surface Chamfer & 20.0 & Surface normals & 16.0 \\
Laplacian & 0.1 & Edge length & 0.1 \\
Normal consistency & 350.0 & Local rigidity & 100.0 \\
Opening position & 10.0 & Opening normal & 0.1 \\
Opening area & 1.0 & Opening overlap & 1.0 \\
Opening boundary smoothness & 0.5 & Sac volume & 0.05 \\
\bottomrule
\end{tabular}
\end{table}

\subsubsection{CVAE training}

We use the same stratified and balanced split as the other models in this
work: 530 training cases, 132 validation cases, and 100 held-out test cases
with seed 42. The CVAE is trained for 4000 epochs with batch size 96 using
AdamW with learning rate $10^{-4}$ and weight decay $10^{-4}$. The latent
dimension is 108, the hidden dimension is 256, the condition embedding
dimension is 64, and the vessel-condition embedding dimension is 96. We use a
KL warm-up of 1000 epochs, free bits of 0.02, validation every 50 epochs, and
select the checkpoint with the lowest validation total loss. The selected
SynVA-A1 checkpoint is from epoch 2550.

The CVAE loss combines three parts. The posterior reconstruction path matches
the fitted GHD coefficients, scale, and decoded vertices while regularizing
the posterior with a KL term. Geometry-aware losses supervise the generated
aneurysm near the ostium through attachment-sample distance, sac-center error,
wrong-side placement relative to the ostium plane, and opening-center error.
Finally, prior-path calibration decodes direct prior samples
$\bm{z}_{\mathrm{prior}}\sim\mathcal{N}(\bm{0},\bm{I})$ during training and
matches them to case-level and batch-level target statistics. We write the
overall objective as
$\mathcal{L}_{\mathrm{CVAE}}=\mathcal{L}_{\mathrm{rec}}+\mathcal{L}_{\mathrm{geo}}+\mathcal{L}_{\mathrm{prior}}$.
The corresponding weights are listed in
Table~\ref{tab:synva-a1-cvae-loss-weights}.

\begin{table}[t]
\centering
\caption{CVAE loss weights used for SynVA-A1.}
\label{tab:synva-a1-cvae-loss-weights}
\scriptsize
\begin{tabular}{lclc}
\toprule
Term & Weight & Term & Weight \\
\midrule
Coeff. MSE & 1.0 & Attachment samples & 5.0 \\
Scale Huber & 1.0 & Sac center & 10.0 \\
Vertex MSE & 250.0 & Side & 2.0 \\
KL & 0.01 & Opening center & 10.0 \\
Prior Huber & 0.25 & Prior mean/std & 0.25 / 8.0 \\
\bottomrule
\end{tabular}
\end{table}

\subsection{Experimental setting}

All ablations are trained and evaluated with the same splits and the same
100-case test set. AneuG-Base denotes the compact GHD-VAE baseline. AneuG-Cond
adds local ostium and vessel conditioning, AneuG-Cond-Morph additionally uses
the morphology vector, and AneuG-Cond-Prior adds prior-path calibration without
morphology conditioning. SynVA-A1 is the selected model with local
conditioning, morphology conditioning, geometry-aware ostium supervision, and
prior-path calibration. SynVA-A1-Ostium uses the same model but increases the
ostium-related supervision weights.

Evaluation uses paired surface Chamfer distance with 10k sampled surface
points, mean absolute errors of morphology descriptors, and distribution-level
metrics in standardized morphology space. Lower values are better for all
reported metrics except coverage.

\subsection{Additional results}

\begin{table}[t]
\centering
\caption{Summary of held-out SynVA-A1 ablation results. Paired metrics are computed against the ground-truth sac; distribution metrics are computed in standardized morphology space. Lower is better except for coverage.}
\label{tab:synva-a1-results-summary}
\scriptsize
\setlength{\tabcolsep}{3.0pt}
\begin{tabular}{lrrrrrrrr}
\toprule
Variant & Chamfer & $D_{\max}$ & $H_{\max}$ & $W_{\max}$ & paired $\bm{z}$-L2 & coverage $\uparrow$ & MMD & Frechet \\
\midrule
AneuG-Base & 0.094099 & 0.113297 & 0.119853 & 0.074605 & 1.284607 & 0.84 & 0.006438 & 1.092620 \\
AneuG-Cond & 0.089148 & 0.106784 & 0.108650 & 0.068531 & 1.217317 & 0.81 & 0.009033 & 1.111498 \\
AneuG-Cond-Morph & \textbf{0.082113} & 0.107940 & 0.103746 & 0.061546 & 1.127876 & 0.83 & 0.010141 & 1.254649 \\
AneuG-Cond-Prior & 0.087192 & \textbf{0.084195} & 0.081362 & 0.059200 & 1.008820 & \textbf{0.87} & \textbf{0.005170} & \textbf{0.984928} \\
SynVA-A1 & 0.086362 & 0.084209 & \textbf{0.079572} & 0.057382 & \textbf{0.980470} & 0.85 & 0.005519 & 1.017597 \\
SynVA-A1-Ostium & 0.086228 & 0.088093 & 0.082618 & \textbf{0.056933} & 0.996746 & 0.83 & 0.006162 & 1.061131 \\
\bottomrule
\end{tabular}
\end{table}

The ablation shows that the individual additions affect different aspects of
the generated aneurysms. Local conditioning improves paired reconstruction,
morphology conditioning gives the lowest Chamfer distance, and prior-path
calibration is most important for distribution-level behavior. SynVA-A1 gives
the best paired morphology distance and the lowest height error, while
AneuG-Cond-Prior achieves the strongest coverage, MMD, and Frechet distance.

\IfFileExists{images/fail_cases_synva_a1.png}{%
\begin{figure}[ht]
  \centering
  \includegraphics[width=0.8\textwidth]{images/fail_cases_synva_a1.png}
  \vspace{-0.2cm}
  \caption{Representative SynVA-A1 failure cases.}
  \label{fig:synva-a1-failure-cases}
\end{figure}
}{}

% Required packages: amsmath, amssymb, bm, booktabs, graphicx

\section{SynVA-A1}
\label{Appendix:SynVA-A1}

\subsection{Model-specific data pre-processing}

SynVA-A1 is trained on the processed real aneurysm cases described in the main
text, i.e. the combined AneuX, IntrA, and CMHA cases after the common mesh and
label preprocessing. We use only those cases for which the GHD fitting step can
be computed successfully. The synthetic SynVA-P1 dataset is not used for this
model. The model-specific preprocessing has two purposes: it first converts
each aneurysm into a common GHD representation, and it then extracts the
condition variables used by the conditional VAE.

For the GHD fit, the target geometry is the aneurysm sac together with the
ostium neck. The surrounding parent vessel and adjacent branches are excluded
from the deformation target, because SynVA-A1 is intended to generate an
aneurysm anchored to a given ostium rather than a complete pathological vessel
segment. The sac mesh remains open at the ostium. The opening is not closed
for canonical-mesh construction; instead, its boundary is retained and later
used to build the ostium information required for the fitting losses.

All fits are expressed relative to one canonical aneurysm mesh. Instead of
choosing one training case as this template, we build an average canonical sac
from the training data. After an initial ostium-based alignment, rays are cast
from a common origin $\bm{c}$ along shared directions $\bm{u}_j$. If $r_{ij}$
is the intersection distance for case $i$ and ray $j$, the canonical radius is
$\bar{r}_j=\operatorname*{median}_{i} r_{ij}$ and the corresponding canonical
vertex is $\bm{v}^{\mathrm{can}}_j=\bm{c}+\bar{r}_j\bm{u}_j$. This produces a
single canonical sac with an open neck boundary and a shared graph topology for
all subsequent fits. The canonical ostium ring is extracted from this boundary.

After the canonical and target meshes have been prepared, an ostium checkpoint
is created for each case. The ostium points are mapped to nearby mesh
vertices, sorted into an ordered boundary ring, optionally smoothed in the
ostium plane, and triangulated as a simple opening patch. This patch is not
used to close the canonical sac; it provides a consistent surface object on
which the neck and opening losses can be evaluated.

Each target aneurysm is then aligned to this canonical frame before
optimization. The pre-alignment is stored as a homogeneous transformation
matrix $\bm{H}_{i}^{\mathrm{pre}}\in\mathbb{R}^{4\times4}$. It translates the
target ostium center to the canonical center, rotates the target ostium normal
onto the canonical normal, and aligns the dominant in-plane axis of the target
ostium with the corresponding canonical axis. For homogeneous coordinates
$\tilde{\bm{p}}$, the aligned point is
$\tilde{\bm{p}}_{i}^{\mathrm{align}}=\bm{H}_{i}^{\mathrm{pre}}\tilde{\bm{p}}_{i}^{\mathrm{raw}}$.
The aligned sac is saved as the fitting mesh, and the stored transformation
allows the fitted result to be related back to the processed data coordinates.
The canonical mesh and the aligned targets are normalized before the GHD
optimization.

The CVAE does not use the raw target mesh as its regression target. Instead,
the target vector is read from the GHD fitting checkpoint and consists of the
flattened GHD coefficients with $K=144$ basis functions and the fitted scale.
The fitted rotation and translation are not predicted by SynVA-A1; they are
only used to decode fitted meshes and to evaluate geometry losses in the
aligned coordinate frame.

The condition input is extracted from the same case and is centered around the
ostium. SynVA-A1 therefore does not use the complete vessel tree as input, but
only the local vessel geometry that is relevant for aneurysm attachment. To
make this vessel patch consistent with the fitted GHD representation, the
candidate vessel points are first mapped into the same local coordinate frame
as the fitted aneurysm. Distances are then measured from the ostium-ring
centroid, and a neighborhood around the ostium is retained. The neighborhood
radius is chosen from the ostium size and the local distance distribution, so
that it covers the neck region and nearby parent-vessel surface without
including the entire vessel. From this local patch, 256 points are selected by
farthest-point sampling. If the local selection contains too few points, the
full candidate surface is used as a fallback.

The remaining case information is prepared in the same local frame. The actual
network condition consists of the ordered ostium ring, compact ostium
parameters containing center, normal, radius, and eccentricity, the local
vessel point set, and a small morphology descriptor vector. We use only
selected global sac descriptors as morphology condition: aneurysm surface area
($A_A$), aneurysm volume ($V_A$), convex-hull surface area ($A_{CH}$),
convex-hull volume ($V_{CH}$), maximum diameter ($D_{\max}$), maximum height
($H_{\max}$), maximum width ($W_{\max}$), orthogonal height ($H_{\perp}$), and
orthogonal width ($W_{\perp}$). In addition, 64 ostium samples and 256 dense
samples from the segmented ostium attachment region are stored for the
geometry losses at the vessel--aneurysm interface. GHD targets and continuous
condition variables are normalized with statistics computed on the training
split.

\subsection{Architecture}

SynVA-A1 consists of a GHD fitting stage and a conditional aneurysm generator.
The fitting stage defines a compact and anatomically aligned parameterization
of each training aneurysm. The generator then learns a conditional prior over
these fitted parameters, using the local ostium, nearby vessel geometry, and
morphology descriptors as conditioning information.

\subsubsection{GHD fitting}

Our fitting procedure is based on the Graph Harmonic Deformation (GHD)
representation introduced in AneuG~\cite{ding2025two}, but adapts it from a
full aneurysm--vessel segment to the local aneurysm domain considered here.
AneuG fits a vascular structure that includes the parent artery and several
openings. In contrast, our target is the aneurysm sac together with its single
ostium neck. The parent vessel is not part of the deformation target and is
used later only as conditioning information for the generator.

We retain the core GHD representation from AneuG. Each case is fitted by
deforming the canonical sac in a shared graph-harmonic basis and by optimizing
a global similarity transform. The target of the fit is therefore not a raw
mesh with arbitrary topology, but a compact set of GHD coefficients in the
common canonical eigenspace. In our implementation this basis contains
$K=144$ graph-harmonic modes.

The general shape terms are also inherited from the original AneuG fitting
objective. We use surface Chamfer distance and surface-normal matching to align
the deformed canonical sac with the target sac. We further keep the standard
mesh regularizers used in AneuG: Laplacian smoothing, edge-length
regularization, normal-consistency regularization, and local rigidity. These
terms control the overall sac shape and discourage noisy or folded
deformations.

The main methodological change is the treatment of the ostium. Original AneuG
supervises several vessel openings and, in its full-vessel setting, can use
differentiable occupancy and differentiable centreline losses to constrain the
parent-vessel segment. Since SynVA-A1 does not fit the parent vessel, these
full-vessel losses are disabled in our sac-and-neck fitting mode. Instead, we
replace the multi-opening objective by a single-ostium objective. The ostium is
represented by an ordered boundary ring and by a reconstructed opening patch
from the ostium checkpoint. This allows us to supervise the neck directly: the
opening position is matched to the target, the opening normal is aligned using
a robust ring-based normal estimate, and an area term prevents collapse of the
neck opening. Additional overlap, boundary-smoothness, and relative-volume
terms stabilize the vessel--aneurysm interface and the global sac size.

Thus, the GHD fitting contribution in SynVA-A1 is not a new deformation
parameterization, but a new anatomical use of the AneuG fitting machinery. We
use an average canonical sac instead of a single representative template, fit
only the aneurysm sac and one open ostium, replace full-vessel centreline and
occupancy constraints by ostium-specific neck constraints, and reuse one
precomputed canonical eigenspace for all target fits.

\subsubsection{Conditional Aneurysm Generator}

SynVA-A1 keeps the compact conditional GHD-VAE backbone used in AneuG. This is
the part that remains deliberately close to the original implementation: the
network operates on GHD parameters rather than mesh vertices, uses a Gaussian
latent space with the standard reparameterization trick, and injects the
condition on the decoder side. In our case the VAE input has 433 dimensions,
corresponding to $144\times3$ GHD coefficients plus the fitted scale.

The encoder receives only this normalized GHD target. It consists of a linear
input projection, one residual MLP block, and two output heads that predict the
latent mean and log-variance. The decoder mirrors this compact design. A
108-dimensional latent sample is concatenated with the encoded condition,
passed through a linear layer and one residual MLP block, and mapped back to
the normalized GHD coefficients and scale. Both encoder and decoder use a
hidden dimension of 256. Thus, the generative backbone is still the same
AneuG-style conditional GHD-VAE: the changes are not in the basic VAE
mechanism, but in what is provided as anatomical condition and how decoded
samples are supervised.

The condition network is the main architectural extension. In the original
AneuG setup, the conditional VAE is driven by a small morphology condition.
SynVA-A1 replaces this by a local vessel--ostium condition. The local vessel
patch contains 256 points around the ostium and is encoded with a
PointNet-style branch using shared point-wise MLP layers followed by max
pooling. The compact ostium parameters, namely center, normal, radius, and
eccentricity, are encoded with a separate MLP. The ordered ostium ring is
resampled to 20 points, flattened, and encoded with another MLP so that the
decoder receives an explicit representation of the neck boundary. These
branches are fused into a 96-dimensional vessel--ostium feature, concatenated
with the selected morphology descriptors, and then projected to a
64-dimensional condition embedding before being passed to the VAE decoder.

SynVA-A1 also differs from the original AneuG model in the training signal
applied to decoded samples. The dense attachment samples from the segmented
ostium region are not used as an additional encoder branch in the final model;
instead, they supervise the generated aneurysm near the vessel--aneurysm
interface. Together with opening-center and side-of-plane losses, this forces
the predicted GHD parameters to respect the prescribed ostium geometry.
Finally, SynVA-A1 adds prior-path calibration: during training, samples drawn
from the standard normal prior are decoded under the same condition and matched
to case-level and batch-level target statistics. This trains the same sampling
path that is used at inference time, where no fitted GHD target is available.

\subsection{Training Details}

\subsubsection{GHD fitting}

GHD fitting is optimized separately for each prepared case with $K=144$ basis
functions and one ostium. We use AdamW with learning rate
$10^{-3}$ for 2000 epochs. The sac-and-neck fitting mode uses one opening and
disables the full-vessel differentiable occupancy and centreline losses from
the original AneuG fitting. The active fitting weights are summarized in
Table~\ref{tab:synva-a1-ghd-loss-weights}.

\begin{table}[t]
\centering
\caption{Loss weights used for the SynVA-A1 GHD fitting.}
\label{tab:synva-a1-ghd-loss-weights}
\scriptsize
\begin{tabular}{lclc}
\toprule
Term & Weight & Term & Weight \\
\midrule
Surface Chamfer & 20.0 & Surface normals & 16.0 \\
Laplacian & 0.1 & Edge length & 0.1 \\
Normal consistency & 350.0 & Local rigidity & 100.0 \\
Opening position & 10.0 & Opening normal & 0.1 \\
Opening area & 1.0 & Opening overlap & 1.0 \\
Opening boundary smoothness & 0.5 & Sac volume & 0.05 \\
\bottomrule
\end{tabular}
\end{table}

\subsubsection{CVAE training}

We use the same stratified and balanced split as the other models in this
work: 530 training cases, 132 validation cases, and 100 held-out test cases
with seed 42. The CVAE is trained for 4000 epochs with batch size 96 using
AdamW with learning rate $10^{-4}$ and weight decay $10^{-4}$. The latent
dimension is 108, the hidden dimension is 256, the condition embedding
dimension is 64, and the vessel-condition embedding dimension is 96. We use a
KL warm-up of 1000 epochs, free bits of 0.02, validation every 50 epochs, and
select the checkpoint with the lowest validation total loss. The selected
SynVA-A1 checkpoint is from epoch 2550.

The CVAE loss combines three parts. The posterior reconstruction path matches
the fitted GHD coefficients, scale, and decoded vertices while regularizing
the posterior with a KL term. Geometry-aware losses supervise the generated
aneurysm near the ostium through attachment-sample distance, sac-center error,
wrong-side placement relative to the ostium plane, and opening-center error.
Finally, prior-path calibration decodes direct prior samples
$\bm{z}_{\mathrm{prior}}\sim\mathcal{N}(\bm{0},\bm{I})$ during training and
matches them to case-level and batch-level target statistics. We write the
overall objective as
$\mathcal{L}_{\mathrm{CVAE}}=\mathcal{L}_{\mathrm{rec}}+\mathcal{L}_{\mathrm{geo}}+\mathcal{L}_{\mathrm{prior}}$.
The corresponding weights are listed in
Table~\ref{tab:synva-a1-cvae-loss-weights}.

\begin{table}[t]
\centering
\caption{CVAE loss weights used for SynVA-A1.}
\label{tab:synva-a1-cvae-loss-weights}
\scriptsize
\begin{tabular}{lclc}
\toprule
Term & Weight & Term & Weight \\
\midrule
Coeff. MSE & 1.0 & Attachment samples & 5.0 \\
Scale Huber & 1.0 & Sac center & 10.0 \\
Vertex MSE & 250.0 & Side & 2.0 \\
KL & 0.01 & Opening center & 10.0 \\
Prior Huber & 0.25 & Prior mean/std & 0.25 / 8.0 \\
\bottomrule
\end{tabular}
\end{table}

\subsection{Experimental setting}

All ablations are trained and evaluated with the same splits and the same
100-case test set. AneuG-Base denotes the compact GHD-VAE baseline. AneuG-Cond
adds local ostium and vessel conditioning, AneuG-Cond-Morph additionally uses
the morphology vector, and AneuG-Cond-Prior adds prior-path calibration without
morphology conditioning. SynVA-A1 is the selected model with local
conditioning, morphology conditioning, geometry-aware ostium supervision, and
prior-path calibration. SynVA-A1-Ostium uses the same model but increases the
ostium-related supervision weights.

Evaluation uses paired surface Chamfer distance with 10k sampled surface
points, mean absolute errors of morphology descriptors, and distribution-level
metrics in standardized morphology space. Lower values are better for all
reported metrics except coverage.

\subsection{Additional results}

\begin{table}[t]
\centering
\caption{Summary of held-out SynVA-A1 ablation results. Paired metrics are computed against the ground-truth sac; distribution metrics are computed in standardized morphology space. Lower is better except for coverage.}
\label{tab:synva-a1-results-summary}
\scriptsize
\setlength{\tabcolsep}{3.0pt}
\begin{tabular}{lrrrrrrrr}
\toprule
Variant & Chamfer & $D_{\max}$ & $H_{\max}$ & $W_{\max}$ & paired $\bm{z}$-L2 & coverage $\uparrow$ & MMD & Frechet \\
\midrule
AneuG-Base & 0.094099 & 0.113297 & 0.119853 & 0.074605 & 1.284607 & 0.84 & 0.006438 & 1.092620 \\
AneuG-Cond & 0.089148 & 0.106784 & 0.108650 & 0.068531 & 1.217317 & 0.81 & 0.009033 & 1.111498 \\
AneuG-Cond-Morph & \textbf{0.082113} & 0.107940 & 0.103746 & 0.061546 & 1.127876 & 0.83 & 0.010141 & 1.254649 \\
AneuG-Cond-Prior & 0.087192 & \textbf{0.084195} & 0.081362 & 0.059200 & 1.008820 & \textbf{0.87} & \textbf{0.005170} & \textbf{0.984928} \\
SynVA-A1 & 0.086362 & 0.084209 & \textbf{0.079572} & 0.057382 & \textbf{0.980470} & 0.85 & 0.005519 & 1.017597 \\
SynVA-A1-Ostium & 0.086228 & 0.088093 & 0.082618 & \textbf{0.056933} & 0.996746 & 0.83 & 0.006162 & 1.061131 \\
\bottomrule
\end{tabular}
\end{table}

The ablation shows that the individual additions affect different aspects of
the generated aneurysms. Local conditioning improves paired reconstruction,
morphology conditioning gives the lowest Chamfer distance, and prior-path
calibration is most important for distribution-level behavior. SynVA-A1 gives
the best paired morphology distance and the lowest height error, while
AneuG-Cond-Prior achieves the strongest coverage, MMD, and Frechet distance.

\IfFileExists{images/fail_cases_synva_a1.png}{%
\begin{figure}[ht]
  \centering
  \includegraphics[width=0.8\textwidth]{images/fail_cases_synva_a1.png}
  \vspace{-0.2cm}
  \caption{Representative SynVA-A1 failure cases.}
  \label{fig:synva-a1-failure-cases}
\end{figure}
}{}
